\documentclass[11pt,a4paper]{article}
\usepackage[utf8]{inputenc}
\usepackage[T1]{fontenc}
\usepackage[margin=1in]{geometry}
\usepackage{hyperref}
\usepackage{setspace}
\setstretch{1.4}

\title{\textbf{DISTANCE}\\ \Large Beyond the Illusion of Spacetime in the Universe}
\author{Kangbin Yim\thanks{Email: yim@sch.ac.kr (Dept. of Information Security Engineering, Soonchunhyang University)}}
\date{May 2025}

\begin{document}

\maketitle
%\newpage

\section*{It Was Never Time - The Universe's True Hardware and the Equalization Engine}
We believe that time flows. We tend to imagine that the events of our lives are arranged along an invisible thread called time, much like gemstones strung one after another on a necklace. The past lies behind us, the future stretches ahead, and every event appears to occupy its own place on that thread, separated from the next by a measurable interval. I wake up in the morning, check the clock, move according to my appointments, and welcome the night as the day ends. The smartphone screen displays time down to the second, and the car dashboard tells me how fast I am moving. We calculate the distance from home to school, from the office to the meeting room, from the Earth to the moon, and from the solar system to other stars. Time and space appear to be the most fundamental realities sustaining our lives.

However, at some point, I began to ask myself this question: Does time truly exist first, and do events occur within it? Does space truly exist first, and do objects move within it? Or are the concepts we call time and space merely names humans attached later to explain a more fundamental mechanism?

I think time, space, distance, and speed are not the hardware of the universe, but closer to an 'interface' created by humans to manage reality. Just as folders and trash bins on a computer screen are symbolic interfaces designed to make complex internal electrical signals easier for humans to handle, time and space might also be a screen constructed for our survival and perception. Our brains evolved not to understand the fundamental structure of the universe, but to predict the trajectory of a thrown stone and judge the distance to a predator.

If so, what is the actual hardware behind this screen of time and space? I believe at its root lies the "Energy Gap". The fundamental cause of a rock rolling down a mountain peak or water flowing down a valley is not time or space, but the difference in potential energy—the energy gap—between the high place and the low ground, and the process of its resolution. From this perspective, the Big Bang can be understood not as an event where matter exploded within empty space, but as the greatest energy gap created by an unimaginable state of energy concentration. For almost an eternity, the universe hasn't simply expanded; it has been endlessly "falling" in the direction of resolving that initial, extreme energy gap.

Here, I also came to rethink the meaning of entropy. We often call entropy chaos or disorder, but I prefer to see it as 'Equalization'. The state where useful energy gaps disappear, flows stop, and we approach a state of complete flatness with no remaining differences is exactly when entropy reaches its maximum.

If the fundamental tendency of the universe is to proceed toward a state where everything becomes flat, why did complex and dynamic structures like life emerge? For a long time, life has been understood as an exceptional existence that draws in negative entropy from the outside to maintain its internal order, resisting the cosmic flow. However, I believe we must take a step further here. The localized order maintained by life is not for resisting the universe; rather, it is merely a structural condition for a far greater increase in entropy.

The process of eating, digesting, and moving is the act of drawing in high-density energy, converting it to low-density energy, and scattering it into our surroundings. On the surface, a living organism seems to maintain the order of its own body, but in reality, it causes a much larger dispersion of energy in the surrounding world to sustain that small order. Viewed this way, life is not a noble exception resisting the universe, but a "Local Engine" that fiercely accelerates the universe's equalization. Every process of eating, digesting, moving, and reproducing is part of the process of resolving energy gaps. I see life not as something outside the laws of the universe, but as a highly advanced energy dispersion structure created by those laws.

\section*{Was the Big Bang Truly an Event of Space Expansion?}
When we talk about the Big Bang, we often picture a balloon. We learn that just as dots on the surface of a small balloon move further apart as it inflates, the universe expanded in the same way. However, at a certain point, I began to question this obvious explanation. Hidden within it is the massive assumption that "space itself actually stretched". Yet, we have never actually seen space stretch. What we observe is only the change in the average density between galaxies, the distribution of energy, and the concentration of matter.

Imagine a transparent cup full of water, into the center of which a single drop of black ink is dropped. At first, an immense concentration of pigment is focused in a very small area. In my language, it is a state where a massive momentum and an extreme energy gap exist. As this change continues, the ink diffuses, the concentration drops, and eventually the entire cup turns a uniform gray. Did the cup grow larger while the ink spread? No. What changed was not the size of the cup, but only the density distribution. If an observer could only see the ink particles without seeing the cup, they might watch the particles move apart and say, "Space is expanding".

I believe the Big Bang is exactly like this. The universe before the Big Bang was a state of unimaginably ultra-high energy concentration. Harboring immense momentum, this state could not be stable, and the energy naturally began to disperse and equalize. We call that the Big Bang.

What if the size of the cup were infinite? Surprisingly, nothing changes. Whether finite or infinite, the ink spreads and the density drops. The essence of the phenomenon does not change at all. I see the history of the universe not as a history of space growing larger, but as a history of concentration dropping and momentum resolving. The galaxies, stars, planets, life, and civilizations we see today are all simply energy islands temporarily clumped together within this massive equalization process. Just as the ink eventually mixes completely, they too will be absorbed into a larger-scale equalization process. It is not because the universe is growing, but because the universe is fiercely mixing itself.

\section*{Straight Lines and Circles Are Walking the Same Path}
Since childhood, we are taught that there are two types of motion: linear motion and circular motion. A car runs in a straight line, and a planet orbits in a circle. Physics also strictly separates the two. But I want to ask if these two are truly fundamentally different phenomena.

We are accustomed to viewing motion from the perspective of an 'independent object'. However, looking at the universe from the perspective of energy organization, completely independent entities do not exist. Atoms, cells, humans, and planets are all "energy islands" holding strong internal binding forces. Therefore, motion is not a change in position, but a change in the 'distance' (energy relationship) between these  islands. The question "What energy relationship is being resolved?" is far more fundamental than "What is moving?".

When momentum from the outside acts on a certain island, if that momentum is transmitted uniformly throughout the organization, the island moves in a single direction. We call this linear motion. On the other hand, if the momentum cannot be transmitted uniformly due to internal binding relationships, a rotation occurs within the organization. In other words, circular motion is not a completely different motion, but a 'failed linear motion' that could not stretch out straight due to strong internal binding forces, appearing instead as an internal circulation.

From this perspective, I also came to redefine the meaning of inertia. Inertia is not a 'will' or 'property' of an object stubbornly trying to maintain its current state. Inertia is 'severance'. It is a state where energy communication with other entities is cut off—a state of absence where no new energy relationship (momentum) can be formed. A stationary rock and a spaceship flying at a constant speed are fundamentally the same.

Ultimately, there are not two types of motion—straight lines and circles—in the universe. There exists only a single process attempting to resolve momentum or distance. If that resolution process proceeds freely without resistance (unbalanced momentum), it becomes a straight line, and if it is trapped in a continuous unbalanced distance (leaning toward a closer island), it becomes a circle. Straight lines and circles are by no means different paths. They are merely walking the same path of cosmic momentum resolution towards different energy islands.

\section*{Gravity, the Non-linear Illusion Created by a Linear Ruler}
If linear and circular motions are fundamentally no different, how should we explain 'gravity' (universal gravitation), where objects with mass attract each other across empty space? I came to view this massive phenomenon from a completely different angle as well.

All energy in the universe is ultimately broken down to its limits, aiming for equalization—an ultimate balance where the mutual momentum between the smallest islands across the universe becomes identical. Within this magnificent process, a powerful directionality arises as local energies attempt to resolve their gap (distance) toward a relatively denser macroscopic energy momentum, pushing to become a 'single island'. I believe gravity is not a separate, independent force, but merely the physical manifestation of this fierce directionality attempting to resolve the distance between macroscopic and local momentums. With only differences in scale, all entities in the universe possess this directionality to resolve massive gaps, and universal gravitation is simply the phenomenon created by this directionality.

Then, how should we understand the 'Curvature of Spacetime' proposed by Einstein? Einstein explained that mass warps spacetime, thereby generating gravity. It is a brilliant insight, but as mentioned earlier, in my worldview, time and space are not actual subjects. The true reality of the universe is a 'non-linear field of momentum' intertwined with the relationships of complex islands.

To observe and measure this massive, churning phenomenon of diffusion and change at our scale, humans inevitably had to invent a 'linear measuring tool'. That is precisely what we classically recognize as time and space.

Imagine overlaying a stiff, linear net (a measuring tool) over a churning, non-linear sea of momentum. Naturally, the net is bound to distort and bend to accommodate the non-linearity. Spacetime itself is not curving. When we attempt to read the dynamic and non-linear momentum field of the universe with a linear tool (spacetime), an inevitable error and distortion occur from the ratio between that linearity and non-linearity. That is precisely what we observe as the 'curvature of spacetime'.

Ultimately, if we set aside the tools of time and space and gaze at the universe's true hardware, there is no magical force acting across empty space, nor is there a self-warping stage. There exists only the dazzling flow of momentum, fiercely resolving distance and attempting to fuse into a single island.

\section*{The Continuum of Life and Death}
We are used to clearly dividing the world into the living and the dead. We believe the living move and reproduce, while the dead do not change. However, if we look at the entire universe as one massive energy flow, this seemingly distinct boundary slowly crumbles.

The universe consists of countless energy islands. Stars, planets, rocks, rivers, and even living beings are all islands holding their own energy. Here, I think what matters is not whether they biologically breathe or not, but 'how much they can interact with each other'.

All cosmic phenomena begin from the 'energy gaps' existing between these islands. Because a gap exists, flows and changes occur, structures are built, and this flow is ultimately a massive equalization process heading toward an average value. The history of the universe is a history of gathering and scattering (Ihapjipsan) where these islands constantly assemble and disassemble. Matter that made up a star becomes a planet, the planet's soil forms a living body, and the living body scatters back into the ocean and soil. We simply cut out a specific section of that process and call it 'life,' and another section 'death'.

From the universe's point of view, being alive and being dead are not different worlds. They are merely differences in the position of 'dynamicity' appearing on a linear energy scale. Life is not a mysterious exception from outside the universe. It is a highly sophisticated momentum structure formed within the massive equalization process, fiercely absorbing, expelling, and exchanging energy with its surroundings.

Death, too, is not a simple end. It is the beginning of a new connection where the matter comprising the living body becomes part of another structure. In this massive flow where the universe infinitely mixes and flattens itself, life and death, creation and extinction never oppose each other. What we call 'life' is also just one scene of the cosmic equalization process where countless energy islands repeatedly gather and scatter, fiercely churning for a fleeting moment.

\section*{The Ecosystem and the Formation of Distance}
Once we understand life as the universe's momentum structure, the ecosystem—where countless living beings are intertwined—approaches us in a completely different light. We often describe the ecosystem as a 'food chain' with a beginning and an end, but I see the ecosystem as a massive network without a beginning or an end. Grass connects to the sun, herbivores to the grass, predators to herbivores, decomposers to predators, and the soil back to the grass. This is not a simple list of individuals, but a massive network where numerous momentum islands are tightly woven to mix energy.

From this perspective, the phenomenon of 'predation' is read differently. We often see a lion hunting a zebra as a cruel struggle for survival. But viewed macroscopically, predation is a fierce act of connection where two momentum structures mix their energy. In this endless collision, the lion and the zebra change each other so that the lion hunts more efficiently and the zebra runs away faster. Predation is not a simple extinction, but the evolutionary driving force that forges new structures and new distances.

In addition to the consumption called predation, life resolves momentum through the combination called 'reproduction'. Early life engaged in simple cloning, but this brought limitations to the exploration space. Thus, life created the 'differentiation of sexes'. Reproduction is not simply a process of increasing numbers. It is a momentum resolution process where two different islands connect to fiercely explore completely new structures and possibilities.

But here the most paradoxical truth of life is revealed. Mating is connection, and predation is connection. This attempt at 'connection,' constantly trying to get closer to mix energy, paradoxically creates a completely new species that can never mate as these continuous changes accumulate—a new 'severance' (distance). Connection births severance, and closeness births distance. Amidst this magnificent contradiction, life is a sophisticated structure of the universe that resolves momentum and endlessly forms new distances.

\section*{The Ecosystem, and the Massive Disconnection Born of Connection}
I do not see life as an exceptional miracle that suddenly fell from outside the universe. The momentum that began in massive stars gradually differentiated into smaller scales, creating highly sophisticated energy structures—that is life. We call this massive network, where countless momentum structures are tightly intertwined, the ecosystem.

We often think of the ecosystem as a one-directional hierarchy called a food chain. And we view the phenomenon of 'predation' within it as a cruel struggle for survival. But if we look from a slight distance, a different landscape appears. A lion hunting a zebra is not mere extinction. It is a fierce act of connection where two momentum structures mix their energy. In this endless collision, the lion and zebra forge each other's structures—the lion to hunt more efficiently, the zebra to run faster.

Beyond the consumption called predation, life also resolves momentum through the combination called 'reproduction'. The differentiation of sexes is not simply a device to increase numbers. It is a process where two different islands connect to fiercely explore new genetic combinations and possibilities.

Yet, in this magnificent flow, a strange paradox arises. Predation is connection, and mating is connection. This attempt at 'connection,' constantly trying to get closer to mix energy, paradoxically creates a new species that can never be crossed as these continuous changes accumulate—a new 'disconnection' or 'distance'. Connection births disconnection, and closeness births distance. In this endless gathering, scattering, and contradiction, life resolves momentum, acting as a precise cogwheel moving steadfastly toward the universe's equalization.

\section*{The New Evolution Chosen by the Most Isolated Predator}
In my childhood, I wondered why humans were so uniquely special. I even imagined that we might be an alien existence that flew in from somewhere in the universe. However, as I looked deeply into the ecosystem, I found a different answer. The more powerful a predator is at the apex of the ecosystem, the fewer similar competitors exist around it. Humanity, too, walked the path of the predator. We drove countless similar hominids, like Neanderthals and Denisovans, to disappear into history.

The price for eliminating all competitors was that humanity fell into the abyss of a massive 'Sexual Distance,' unable to mix genes with any other species on Earth. Pushed off the continent of evolution, we became a thoroughly isolated 'Island Species.'

However, our drive as an equalization engine trying to mix the universe's energy did not stop there. When the traditional path of evolving the biological body was blocked, humanity began to create entirely new entities outside of its own body. Language emerged, writing was born, cities and industries were built, and finally, computers and artificial intelligence were created. We completely shifted our path from the combination of genes to the combination of civilizations, from the evolution of the individual to the evolution of massive systems.

When the door to biological exchange closed, humanity, the most powerful predator, became the existence that creates the largest and fiercest 'external equalization engines' in the universe.

\section*{The Equalization Engine of Society and the Dispersion of Power}
We learn human history as a process of great moral progress where freedom and equality have expanded. Monarchies collapsed, slavery was abolished, and democracy took root. However, from my perspective, history is not a conflict between good and evil. It is a physical equalization process where the momentums formed within the massive energy structure of society accumulate, shatter, and resolve.

For survival, early humans concentrated decision-making power in a few kings and nobles. Power was precisely the 'structural node' where the momentum of the entire society gathered. But this powerful concentration formed an extreme 'energy gap' between the rulers and the ruled. The most brutal manifestation of this energy gap was slavery.

When an energy gap condenses, a massive momentum inevitably boils up. Citizen revolutions, like the French Revolution, were not justice dropping out of the sky. They were explosions of momentum resolution, where the fiercely accumulated social potential over hundreds of years poured violently toward a new structure.

The democracy and fairness we pride ourselves on are also highly advanced 'momentum dispersion technologies' that finely distribute a massive, clustered power gap to tens of millions of individual citizens. The movement from a clustered, low-resolution society to a finely fragmented, high-resolution, homogeneous society. The history written in human blood and sweat might ultimately be a magnificent thermodynamic trajectory in which our society has fiercely shaved down its own gaps toward the ultimate flattening of the universe.

\section*{The Last Asymmetry Remaining at the End of Fair Competition}
Civilization has relentlessly reduced asymmetries. The differences in social status, political opportunities, and educational gaps have been flattened out. We cheer this as the expansion of fairness, but fairness never extinguishes competition. Rather, it draws more people into the same arena, 'universalizing competition.'

This trajectory of equalization appears most dramatically in the relationship between men and women. In the past, due to survival environments centered on physical strength and manual labor, men and women lived on different tracks. However, the development of technology and industry erased the importance of muscle power and merged the two tracks into one giant arena. Now, men and women compete fiercely under the exact same system, in the same classrooms and the same workplaces.

In this empty space where all the asymmetries of society have been smoothly shaved away, only one massive biological asymmetry, which human institutions can never shave down, stands tall alone. It is 'pregnancy and childbirth,' which occurs exclusively through the female body.

Today's low birth rate is not a phenomenon caused simply by a lack of money or selfishness. It is the rupturing sound created by a head-on collision between the 'civilizational trajectory of equalization'—which eliminates all differences to perfect symmetric fair competition—and the 'biological structure of reproduction,' which is inherently bound to be asymmetric. The reason birth rates plummet as countries become more advanced is not because they are the wealthiest, but because they are the 'most symmetric and flat societies' where men and women run on the exact same track.

\section*{Cosmic Gathering and Scattering, and the Dash of Information Civilization}
Looking at my reflection in the mirror, I always believe I am the same entity, but the matter making up my body constantly escapes and is replaced while I breathe and eat. Life is not a fixed lump of matter, but a dynamic 'flow' where energy fiercely enters and exits. Being alive and being dead are not two different worlds. When energy gathers and burns fiercely, it becomes life; when that connection breaks and hardens, it becomes death. Death is not an end, but a magnificent gathering and scattering (Ihapjipsan) returning to cosmic dust to be reorganized into new flows.

And humanity has transcended this gathering and scattering of the finite physical body. A hunter's memory carved into a mural becomes writing and books, transcending physical limits to bind a human a thousand years later with a human a thousand kilometers away into one. At the end of this massive 'Long Distance' network lies modern information civilization and the Internet.

Just because information is invisible does not mean it has escaped the thermodynamic shackles of the universe. The Internet and Artificial Intelligence (AI) swallow massive amounts of electricity and emit enormous heat. We sit still, scrolling through smartphones, but in reality, our thoughts and data are breaking down and mixing the universe's energy more explosively than ever before. This brilliant civilization we have built is not a fortress of order, but the ultimate massive accelerator the universe operates to erase its own differences.

\section*{(Epilogue) We Are Passing Through Energy Gaps}
Standing at the end of this long contemplation, a chilling question inevitably arises. If both life and civilization are merely mixers for mixing and dispersing the universe's energy, what meaning is there in human thought and emotion, art and love?

I actually draw deep comfort from this fact. We are not isolated dust randomly thrown into cosmic space. Every moment I breathe, eat, suffer, and create is a cogwheel perfectly connected to the universe's grand equalization process. To love someone is to build a bridge between two disconnected energy islands, creating a completely new flow and a cosmic event.

Of course, the fact that humans are such powerful accelerator pedals leaves us with a heavy homework. Grinding up hundreds of millions of years of fossil fuels buried in the Earth's strata in just a few centuries and mixing energy at a terrifying speed, how we will responsibly handle this cosmic flow is now the greatest question humanity faces.

Now, when the second hand of a clock moves, I do not simply believe time flows. I see the energy fiercely resolving behind it. We are not living time. We are passing through energy gaps. We are not fixed entities sitting alone in space, but dazzling flows temporarily formed within the field of relationships. Until that distant destination where all differences vanish, we will remain as vortexes dancing the most fierce and beautiful dance as a fleeting burst in the middle of the cosmic flow.

\end{document}